%
% Teil 4: Der Stab endlicher Laenge -- keine Verzweigung, sondern
% unendlich viele Pole auf der negativen reellen Achse;
% Residuenreihe (Fourier-Reihe) und Grenzuebergang l -> unendlich
%
\section{Der Stab endlicher Länge: unendlich viele Pole statt
Verzweigungsschnitt}
\label{laplace:subsection:endlicher_stab}

Wir betrachten nun denselben Wärmeleiter, aber mit endlicher Länge $l$.
Am linken Ende $x = 0$ wird die Temperatur $A_0(t)$ eingeprägt, das rechte
Ende $x = l$ wird auf null gehalten, und die Anfangstemperatur verschwindet.
Die Bildgleichung~\eqref{laplace:eq:bildgleichung} bleibt unverändert, nur
treten an die Stelle der Beschränktheitsforderung jetzt zwei
Randbedingungen.
Wie beim halbunendlichen Stab spalten wir die Lösung in die Randanregung
$a_0(s)$ und eine Basislösung $u_0$ auf, die am linken Rand den Wert eins
und am rechten den Wert null annimmt.
Die Basislösung gewinnen wir aus der allgemeinen exponentiellen Lösung
$c_1 e^{x\sqrt{s}} + c_2 e^{-x\sqrt{s}}$ von
Gleichung~\eqref{laplace:eq:bildgleichung}:
Die Randbedingung bei $x = l$ erzwingt die antisymmetrische Kombination
$\sinh\bigl((l-x)\sqrt{s}\bigr)
= \tfrac{1}{2}\bigl(e^{(l-x)\sqrt{s}} - e^{-(l-x)\sqrt{s}}\bigr)$,
die dort verschwindet, und die Normierung auf den Wert eins bei $x = 0$
liefert den Nenner $\sinh\bigl(l\sqrt{s}\bigr)$.
So ergibt sich
\begin{equation}
    u(x,s)
    =
    a_0(s) \, u_0(x,s),
    \qquad
    u_0(x,s)
    =
    \frac{\sinh\bigl((l-x)\sqrt{s}\bigr)}{\sinh\bigl(l\sqrt{s}\bigr)},
    \label{laplace:eq:stab_bildfunktion}
\end{equation}
wie ebenfalls in \cite{laplace:doetsch1961} vorgerechnet wird.
Um die Formeln übersichtlich zu halten, setzen wir im Folgenden $l = \pi$;
der allgemeine Fall entsteht daraus durch einfache Skalierung.

Wieder steht $\sqrt{s}$ in der Bildfunktion, und nach dem letzten Abschnitt
erwarten wir reflexartig einen Verzweigungsschnitt.
Doch die Erwartung trügt.
Ersetzen wir probeweise $\sqrt{s}$ durch $-\sqrt{s}$, was genau dem
Wechsel auf den anderen Zweig der Wurzel entspricht, so wechseln wegen
$\sinh(-z) = -\sinh(z)$ Zähler und Nenner \emph{gleichzeitig} das
Vorzeichen:
\begin{equation}
    \frac{\sinh\bigl(-(\pi-x)\sqrt{s}\bigr)}
         {\sinh\bigl(-\pi\sqrt{s}\bigr)}
    =
    \frac{-\sinh\bigl((\pi-x)\sqrt{s}\bigr)}
         {-\sinh\bigl(\pi\sqrt{s}\bigr)}
    =
    u_0(x,s).
    \label{laplace:eq:eindeutigkeit}
\end{equation}
Die Funktion $u_0$ merkt vom Zweigwechsel überhaupt nichts, sie ist
\emph{eindeutig}.
Noch deutlicher wird das in der Potenzreihe, denn $\sinh$ ist ungerade, und
deshalb enthält
\begin{equation}
    \sinh\bigl(a\sqrt{s}\bigr)
    =
    \sqrt{s}
    \left(
        a + \frac{a^3}{3!}\,s + \frac{a^5}{5!}\,s^2 + \cdots
    \right)
    \label{laplace:eq:sinh_reihe}
\end{equation}
die Wurzel nur als gemeinsamen Vorfaktor, der sich im Quotienten
$u_0(x,s)$ herauskürzt.
Übrig bleibt eine Potenzreihe in $s$ selbst.
Das ist der angekündigte Aha-Moment:
Die endliche Geometrie hat den Verzweigungsschnitt zum Verschwinden
gebracht, und $u_0$ ist eine meromorphe Funktion, deren einzige
Singularitäten isolierte Pole sind.

Diese Pole sind die Nullstellen des Nenners.
Der komplexe Sinus Hyperbolicus hängt über die Identität
$\sinh(z) = -i \sin(iz)$ mit dem Sinus zusammen; seine Nullstellen liegen
deshalb sämtlich auf der imaginären Achse, bei $z = i\nu\pi$ mit
ganzzahligem $\nu$.
Aus $\sinh(\pi\sqrt{s}) = 0$ folgt somit $\pi\sqrt{s} = i\nu\pi$, also
\begin{equation}
    s_\nu
    =
    -\nu^2,
    \qquad
    \nu = 1, 2, 3, \dots,
    \label{laplace:eq:polstellen}
\end{equation}
wobei $s = 0$ ausscheidet, weil dort auch der Zähler verschwindet und die
Singularität hebbar ist.
Die Bildfunktion besitzt somit \emph{unendlich viele} einfache Pole, die
sich mit quadratisch wachsenden Abständen auf der negativen reellen Achse
aufreihen, wie Abbildung~\ref{laplace:fig:pol_kette} zeigt.

\begin{figure}
    \centering
    % ============================================================================
% Bild 4: Bemasste "technische Zeichnung" der unendlichen Pol-Kette
% ============================================================================
% MATHEMATISCHER HINTERGRUND / DESIGN-ENTSCHEIDUNGEN:
%
% * Dargestellt werden die Pole s_nu = -nu^2 (nu = 1, 2, 3, ...) der
%   Bildfunktion u_0(x,s) = sinh((pi-x)sqrt(s))/sinh(pi sqrt(s)) des
%   Stabes der Laenge l = pi, sowie die Bogenfolge C_n aus dem Paper
%   mit den im URSPRUNG zentrierten Radien R_n = (n + 1/2)^2, die
%   GENAU ZWISCHEN den Polen hindurchlaufen:
%       R_1 = 2.25   (zwischen -1 und -4)
%       R_2 = 6.25   (zwischen -4 und -9)
%
% * MASSSTAB: 0.4 cm pro s-Einheit (linear bis zur Bruchstelle).
%   Damit liegen die Pole bei x = -0.4, -1.6, -3.6 cm und die
%   Bogenradien betragen 0.9 cm bzw. 2.5 cm.
%
% * BOGEN-GEOMETRIE: Die Boegen schliessen an die Bromwich-Gerade
%   Re(s) = gamma an (gamma = 1 s-Einheit = 0.4 cm), d.h. die Gerade
%   ist SEHNE der Kreise (konsistent zur D-Kontur-Abbildung).
%   Startwinkel aus cos(theta) = 0.4/R:
%       C_1: R = 0.9  -> theta = 63.61 deg, Endpunkt-y = 0.9*sin = 0.806
%       C_2: R = 2.5  -> theta = 80.79 deg, Endpunkt-y = 2.5*sin = 2.468
%
% * BEMASSUNG: Massketten unterhalb der Achse zeigen die Abstaende
%   1, 3, 5 (in s-Einheiten) zwischen 0,-1,-4,-9 -- die ungeraden
%   Zahlen machen das quadratische Wachstum |s_nu| = nu^2 sichtbar.
%   Technischer Stil: duenne graue Hilfslinien + Masslinien mit
%   Endstrichen (Bar-Tips statt Pfeilen, weil das kuerzeste Mass mit
%   0.4 cm zu schmal fuer beidseitige Pfeilspitzen ist).
%
% * BRUCHSTELLE: Doppelter Zickzack-Bruch (techn. Zeichnungsnorm)
%   nach -9; dahinter ein letzter Pol "-n^2" und ein Achsenstummel
%   mit Pfeil nach links (Fortsetzung der Kette bis -unendlich).
%
% * Bei s = 0 KEIN Pol (hebbare Singularitaet, Zaehler verschwindet
%   mit): offener Kreis statt Kreuz, graue Annotation "hebbar".
% ============================================================================
\scalebox{0.85}{
\begin{tikzpicture}

% ==============================
% 1. FARBEN DEFINIEREN
% ==============================
\definecolor{plotblue}{RGB}{0, 0, 200}
\definecolor{plotred}{RGB}{220, 30, 30}
\definecolor{plotgray}{RGB}{140, 140, 140}
\definecolor{plotbg}{RGB}{244, 244, 252}

% ==============================
% 2. BOEGEN C_n (zuunterst, damit alles andere darueber liegt)
% ==============================
% C_1: R = 0.9 cm, Sehne bei x = 0.4  ->  arc von 63.61 bis 296.39 Grad
\draw[thick, plotgray, dashed] (0.4, 0.806) arc (63.61:296.39:0.9);
% C_2: R = 2.5 cm, Sehne bei x = 0.4  ->  arc von 80.79 bis 279.21 Grad
\draw[thick, plotgray, dashed] (0.4, 2.468) arc (80.79:279.21:2.5);

% Bogen-Labels (jeweils knapp ausserhalb des Bogens, oben links)
\node[plotgray, font=\normalsize] at (-0.55, 1.15) {$C_1$};
\node[plotgray, font=\normalsize] at (-1.75, 2.15) {$C_2$};

% ==============================
% 3. BROMWICH-GERADE (Sehne der Boegen bei Re(s) = gamma)
% ==============================
\draw[thick, plotblue, decoration={markings, mark=at position 0.6 with {\arrow{Stealth[length=3mm, width=2mm]}}}, postaction={decorate}]
    (0.4, -2.468) -- (0.4, 2.468);
\draw[thick] (0.4, 0.1) -- (0.4, -0.1) node[below right, inner sep=2pt] {$\boldsymbol{\gamma}$};

% ==============================
% 4. ACHSEN MIT BRUCHSTELLE
% ==============================
% Hauptachse (rechts der Bruchstelle)
\draw[thick, -{Stealth[length=3.5mm, width=2.5mm]}] (-3.9, 0) -- (2.3, 0) node[right] {\textbf{Re}$(s)$};
% Achsenstummel links der Bruchstelle, Pfeil nach links (bis -unendlich)
\draw[thick, -{Stealth[length=3.5mm, width=2.5mm]}] (-4.34, 0) -- (-6.05, 0);
% Doppelter Zickzack-Bruch (techn. Norm fuer "hier fehlt ein Stueck")
\draw[thick] (-3.92, -0.3) -- (-4.04, -0.06) -- (-3.96, 0.06) -- (-4.08, 0.3);
\draw[thick] (-4.14, -0.3) -- (-4.26, -0.06) -- (-4.18, 0.06) -- (-4.30, 0.3);
% Kurze Im-Achse
\draw[thick, -{Stealth[length=3.5mm, width=2.5mm]}] (0, -2.9) -- (0, 3.05) node[above] {\textbf{Im}$(s)$};

% ==============================
% 5. POLE (Kreuze) UND HEBBARE STELLE
% ==============================
\newcommand{\pole}[3][plotred]{
    \draw[thick, #1, line cap=round] (#2-0.12,#3-0.12) -- (#2+0.12,#3+0.12);
    \draw[thick, #1, line cap=round] (#2-0.12,#3+0.12) -- (#2+0.12,#3-0.12);
}

% Pole bei s = -1, -4, -9 (Massstab 0.4 cm/Einheit) und -n^2
\pole{-0.4}{0}
\pole{-1.6}{0}
\pole{-3.6}{0}
\pole{-5.0}{0}

% Pol-Beschriftungen OBERHALB der Achse (unten laeuft die Masskette)
\node[anchor=south, font=\small] at (-0.4, 0.22) {$-1$};
\node[anchor=south, font=\small] at (-1.6, 0.22) {$-4$};
\node[anchor=south, font=\small] at (-3.6, 0.22) {$-9$};
\node[anchor=south, font=\small] at (-5.0, 0.22) {$-n^2$};

% Hebbare Singularitaet bei s = 0: roter Kreis mit rotem Kreuz
\draw[thick, plotred] (0,0) circle (3.0pt);
\draw[thick, plotred, line cap=round] (-0.075,-0.075) -- (0.075,0.075);
\draw[thick, plotred, line cap=round] (-0.075,0.075) -- (0.075,-0.075);



% ==============================
% 7. LEGENDE
% ==============================
\node[draw, thick, rounded corners=3pt, fill=white, inner xsep=5pt, inner ysep=4pt, anchor=north west] at (1.3, 3.05) {
    \renewcommand{\arraystretch}{1.2}
    \begin{tabular}{@{}c l@{}}

        \tikz[baseline=-0.6ex]{
            \draw[thick, plotred, line cap=round] (-0.1,-0.1) -- (0.1,0.1) (-0.1,0.1) -- (0.1,-0.1);
        } & Polstelle \\

        \tikz[baseline=-0.6ex]{
            \draw[thick, plotred] (0,0) circle (3.0pt);
            \draw[thick, plotred, line cap=round] (-0.075,-0.075) -- (0.075,0.075);
            \draw[thick, plotred, line cap=round] (-0.075,0.075) -- (0.075,-0.075);
        } & Hebbare Singularit\"at \\

        \tikz[baseline=-0.6ex]{
            \draw[thick, plotblue, ->-=0.6] (-0.3, 0) -- (0.3, 0);
        } & Bromwich-Pfad \\

        \tikz[baseline=-0.6ex]{
            \draw[thick, plotgray, dashed] (-0.3, 0) -- (0.3, 0);
        } & Bogen $C_n$ \\

    \end{tabular}
};

\end{tikzpicture}
}

    \caption{Die Singularitäten des Stabes endlicher Länge:
    unendlich viele einfache Pole bei $s_\nu = -\nu^2$ auf der negativen
    reellen Achse.
    Die Abstände wachsen quadratisch, die Wellenlinie deutet die
    Fortsetzung der Kette bis zum $n$-ten Pol und darüber hinaus an.
    Die gestrichelten Bögen $C_n$ verlaufen jeweils zwischen zwei Polen
    und schliessen den Bromwich-Pfad für die Residuenauswertung.}
    \label{laplace:fig:pol_kette}
\end{figure}

Die Residuen an diesen Polen berechnen wir mit der Quotientenregel:
Hat $f_2$ bei $s_0$ eine einfache Nullstelle und ist $f_1(s_0) \neq 0$, so
ist das Residuum von $\frac{f_1}{f_2}$ dort gleich
$\frac{f_1(s_0)}{f_2'(s_0)}$.
Mit $f_2(s) = \sinh(\pi\sqrt{s})$ und
$f_2'(s) = \frac{\pi}{2\sqrt{s}}\cosh(\pi\sqrt{s})$ ergibt sich nach
kurzer Rechnung mit $\sqrt{s_\nu} = i\nu$ das erfreulich einfache Resultat
\begin{equation}
    \operatorname{Res}\bigl(u_0(x,s), -\nu^2\bigr)
    =
    \frac{2}{\pi} \, \nu \sin(\nu x).
    \label{laplace:eq:stab_residuen}
\end{equation}

Für die Inversion können wir die D-Kontur nicht mehr in einem Zug
schliessen, denn jeder feste Kreisbogen würde nur endlich viele der
unendlich vielen Pole erfassen, und ein Bogen durch einen Pol hindurch ist
verboten.
Stattdessen schliessen wir den Bromwich-Pfad durch eine \emph{Folge} von
wiederum im Ursprung zentrierten Bögen $C_n$, deren Radien
\begin{equation}
    R_n
    =
    \Bigl(n + \tfrac{1}{2}\Bigr)^2
    \label{laplace:eq:bogenradien}
\end{equation}
zwischen den Polen $-n^2$ und $-(n+1)^2$ hindurchlaufen.
Exakt mittig liegt der Radius dabei nur im transformierten
$\sqrt{s}$-Raum, wo $\sqrt{R_n} = n + \tfrac{1}{2}$ die Nullstellen $n$
und $n+1$ halbiert; in der $s$-Ebene selbst liegt
$R_n = n^2 + n + \tfrac{1}{4}$ etwas näher am inneren Pol, was für das
Argument aber keine Rolle spielt.
Der $n$-te Bogen schliesst die ersten $n$ Pole ein:
Für jedes feste $n$ liefert der Residuensatz auf der geschlossenen Kurve
aus Geradenstück und Bogen $C_n$ die endliche Partialsumme der ersten $n$
Residuen.
Es wird also nichts von den Bögen \emph{subtrahiert} -- ihr Beitrag geht
für $n \to \infty$ schlicht gegen null, und die Partialsummen konvergieren
gegen die unendliche Residuenreihe.
Für die Anregung mit einem Wärmeimpuls $A_0(t) = \delta(t)$, also
$a_0(s) = 1$, erhalten wir so
\begin{equation}
    U_0(x,t)
    =
    \lim_{n \to \infty}
    \frac{2}{\pi}
    \sum_{\nu=1}^{n}
    \nu \sin(\nu x) \, e^{-\nu^2 t}
    =
    \frac{2}{\pi}
    \sum_{\nu=1}^{\infty}
    \nu \sin(\nu x) \, e^{-\nu^2 t}.
    \label{laplace:eq:stab_loesung}
\end{equation}
Ein Wort zur Sorgfalt ist hier angebracht.
Der Nachweis, dass die Bogenintegrale tatsächlich verschwinden, ist der
entscheidende und in technischen Texten am häufigsten unterschlagene
Schritt; ohne ihn hat die Reihe nur heuristischen Charakter, wie
\cite{laplace:doetsch1961} mit Nachdruck betont.
Das Lemma von Jordan aus
Abschnitt~\ref{laplace:subsection:residuensatz} greift hier nämlich nicht
unmittelbar:
Es verlangt, dass die Bildfunktion für $|s| \to \infty$
\emph{gleichmässig} gegen null strebt, doch wegen der unendlich vielen
Pole kommt der Nenner $\sinh(\pi\sqrt{s})$ in jeder noch so weit
draussen liegenden Umgebung der negativen reellen Achse seinen
Nullstellen beliebig nahe -- von gleichmässigem Abfall kann keine Rede
sein.
Der Nachweis muss deshalb gezielt auf den Radien $R_n$ geführt werden,
die den Polen ausweichen und auf denen $|\sinh(\pi\sqrt{s})|$ nach unten
beschränkt bleibt.
Für die hier auftretenden Bildfunktionen lässt sich der Nachweis führen,
und die entsprechenden Abschätzungen auf den grossen Kreisbögen werden in
\cite{laplace:carslaw1959} für eine ganze Klasse von
Wärmeleitungsproblemen sauber ausgearbeitet.

Das Ergebnis~\eqref{laplace:eq:stab_loesung} verdient einen zweiten Blick,
denn wir kennen es bereits aus einem ganz anderen Kontext.
Es ist exakt die Fourier-Reihe, die die klassische Separation der
Variablen liefert:
Eigenmoden $\sin(\nu x)$ des endlichen Stabes, die mit den Raten
$e^{-\nu^2 t}$ abklingen.
Die Pole der Bildfunktion sind nichts anderes als die Eigenwerte des
räumlichen Problems, und die Residuen sind die zugehörigen Eigenmoden.
Der Residuensatz hat die Modalzerlegung nicht vorausgesetzt, sondern
selbstständig wiederentdeckt.

Damit klärt sich auch der Unterschied zum halbunendlichen Stab endgültig.
Für allgemeine Länge $l$ liegen die Pole bei
$s_\nu = -\frac{\nu^2\pi^2}{l^2}$, ihr Abstand ist also von der Ordnung
$\frac{\pi^2}{l^2}$.
Lassen wir $l$ wachsen, so rücken die Pole immer dichter zusammen, und im
Grenzfall $l \to \infty$ verschmilzt die diskrete Polkette zu einem
Kontinuum auf der negativen reellen Achse, eben dem Verzweigungsschnitt
aus Abschnitt~\ref{laplace:subsection:verzweigung}.
Gleichzeitig geht die Residuenreihe über die Raten $e^{-\nu^2\pi^2 t/l^2}$
in das Schnittintegral über die Raten $e^{-rt}$ aus
Gleichung~\eqref{laplace:eq:keyhole_loesung} über.
Diskretes Spektrum bei endlicher Geometrie, kontinuierliches Spektrum bei
unendlicher Geometrie:
Die komplexe $s$-Ebene bildet die Physik des Problems mit bestechender
Klarheit ab.
